% =====================================================================
%  CAPITULO 1 - continuacao
%  1.4 Potencial eletrico e superficies equipotenciais
% =====================================================================
\subsection{Potencial elétrico e superfícies equipotenciais}

\begin{conceito}[Antes de começar]
\textbf{Potencial de uma carga puntiforme:}
\[
V = \ke\,\frac{q}{d}\quad(\si{\volt}),
\]
com o \emph{sinal} da carga (positivo para $q>0$, negativo para $q<0$).\par
\textbf{Diferença de potencial (ddp)} e trabalho: o trabalho para levar uma carga
$q$ entre dois pontos é $W = q\,U_{AB} = q\,(V_A - V_B)$.\par
\textbf{Campo uniforme:} $U = E\,d$, onde $d$ é a distância medida \emph{ao longo}
da direção do campo.\par
\textbf{Superfície equipotencial:} conjunto de pontos com o mesmo potencial. O
campo elétrico é sempre \textbf{perpendicular} às equipotenciais e aponta no
sentido dos potenciais \emph{decrescentes}. Não há trabalho ao mover uma carga
sobre uma mesma equipotencial.
\end{conceito}

\legendaniveis

\blocofixacao

\begin{questao}[1]{}
Sobre superfícies equipotenciais, é correto afirmar que:
\begin{alternativas}
\item o campo elétrico é tangente a elas.
\item o campo elétrico é perpendicular a elas.
\item o potencial varia ao longo de uma mesma superfície.
\item é necessário realizar trabalho para mover uma carga sobre a superfície.
\item elas sempre coincidem com as linhas de campo.
\end{alternativas}
\resposta{(b)}
\resolucao{Por definição, o potencial é constante sobre uma equipotencial (isso
elimina c). Como não há variação de potencial ao longo dela, o trabalho para mover
uma carga sobre a superfície é nulo (elimina d). O campo, que aponta na direção de
maior variação do potencial, é sempre \emph{perpendicular} às equipotenciais ---
nunca tangente, o que elimina (a) e (e).}
\end{questao}

\begin{questao}[1]{}
Considere dois pontos A e B de um campo elétrico uniforme de intensidade
$E = \SI{e4}{\newton\per\coulomb}$, separados por $\SI{5}{\centi\meter}$ ao longo
da direção do campo. Calcule a ddp entre A e B.
\resposta{$U_{AB} = \SI{500}{\volt}$}
\resolucao{Em campo uniforme, $U = E\,d = \pot{1}{4}\cdot0{,}05 = \SI{500}{\volt}$
(com A no potencial maior, a favor do campo).}
\end{questao}

\blocoaplicacao

\begin{questao}[2]{}
A figura representa três superfícies equipotenciais de um campo elétrico
uniforme, com $V_A = \SI{60}{\volt}$, $V_B = \SI{40}{\volt}$ e
$V_C = \SI{20}{\volt}$, igualmente espaçadas de $\SI{2}{\centi\meter}$.

\circuitoq{%
\begin{tikzpicture}[scale=1,line width=0.8pt,>=Stealth]
 \foreach \x/\v in {0/A,2/B,4/C}{
   \draw[azulprof,thick] (\x,-1.3) -- (\x,1.3);
   \node[above,azulprof] at (\x,1.3) {$V_{\v}$};}
 \node[below] at (0,-1.3) {\footnotesize$\SI{60}{\volt}$};
 \node[below] at (2,-1.3) {\footnotesize$\SI{40}{\volt}$};
 \node[below] at (4,-1.3) {\footnotesize$\SI{20}{\volt}$};
 % campo (do maior para o menor potencial)
 \foreach \y in {-0.7,0,0.7}
   \draw[->,Vcol,thick] (0.3,\y) -- (3.7,\y);
 \node[Vcol] at (2,1.0) {$\vec{E}$};
 \draw[<->,cinza] (0,-1.9) -- (2,-1.9) node[midway,below]{\footnotesize$\SI{2}{\centi\meter}$};
\end{tikzpicture}}{Superfícies equipotenciais em campo uniforme.}

\begin{itensq}
\item Determine a intensidade do campo elétrico.
\item Qual é o trabalho para levar uma carga de $\SI{5}{\micro\coulomb}$ de A até C?
\end{itensq}
\resposta{(a) $E = \SI{1000}{\volt\per\meter}$; (b) $W = \SI{2e-4}{\joule}$}
\resolucao{(a) Entre superfícies adjacentes, $U = \SI{20}{\volt}$ em
$d = \SI{2}{\centi\meter}$:
\[
E=\frac{U}{d}=\frac{20}{0{,}02}=\SI{1000}{\volt\per\meter}.
\]
(b) $U_{AC} = V_A - V_C = 60 - 20 = \SI{40}{\volt}$:
\[
W = q\,U_{AC} = \pot{5}{-6}\cdot40 = \pot{2}{-4}\,\si{\joule}=\SI{2e-4}{\joule}.
\]
O trabalho é positivo porque a carga positiva se move no sentido do campo
(de maior para menor potencial).}
\end{questao}

\begin{questao}[2]{}
A figura mostra linhas equipotenciais no campo gerado por duas cargas de mesmo
módulo e sinais opostos. São dados $V_A = \SI{30}{\volt}$, $V_B = \SI{10}{\volt}$
e $V_C = \SI{-10}{\volt}$. Calcule a ddp entre A e B, e entre B e C.

\circuitoq{%
\begin{tikzpicture}[scale=0.85,line width=0.8pt]
 % duas cargas
 \shade[ball color=Vcol!25] (-3,0) circle (0.3); \node at (-3,0){\footnotesize$+$};
 \shade[ball color=Icol!25] (3,0) circle (0.3);  \node at (3,0){\footnotesize$-$};
 % equipotenciais (arcos estilizados)
 \foreach \r/\lab/\x in {1.2/A/-1.6, 2.0/B/0, 1.2/C/1.6}{}
 \draw[azulprof,thick] (-1.6,-1.6) .. controls (-2.2,0) .. (-1.6,1.6);
 \node[azulprof] at (-1.6,1.85) {$A$};
 \draw[azulprof,thick] (0,-1.9) -- (0,1.9);
 \node[azulprof] at (0,2.15) {$B$};
 \draw[azulprof,thick] (1.6,-1.6) .. controls (2.2,0) .. (1.6,1.6);
 \node[azulprof] at (1.6,1.85) {$C$};
 \node[below,font=\footnotesize] at (-1.6,-1.7) {$\SI{30}{\volt}$};
 \node[below,font=\footnotesize] at (0,-2.0) {$\SI{10}{\volt}$};
 \node[below,font=\footnotesize] at (1.6,-1.7) {$\SI{-10}{\volt}$};
\end{tikzpicture}}{Linhas equipotenciais de um dipolo elétrico.}

\resposta{$U_{AB} = \SI{20}{\volt}$; $U_{BC} = \SI{20}{\volt}$}
\resolucao{A ddp é sempre a diferença dos potenciais:
\[
U_{AB}=V_A-V_B=30-10=\SI{20}{\volt},
\qquad
U_{BC}=V_B-V_C=10-(-10)=\SI{20}{\volt}.
\]
Embora $U_{AB}=U_{BC}$, note que $V_C$ é \emph{negativo}: perto da carga negativa
o potencial fica abaixo de zero. O ponto médio B, equidistante das duas cargas de
módulos iguais, tem potencial intermediário.}
\end{questao}

\begin{questao}[2]{}
As linhas cheias da figura representam linhas de força de um campo eletrostático
e as tracejadas, linhas equipotenciais, com $V_A > V_B > V_C$. Uma partícula de
carga $q = \SI{2e-6}{\coulomb}$ é levada de A até C.

\circuitoq{%
\begin{tikzpicture}[scale=1,line width=0.8pt,>=Stealth]
 % linhas equipotenciais (verticais tracejadas)
 \foreach \x/\lab/\v in {0/A/{V_A},2.2/B/{V_B},4.4/C/{V_C}}{
   \draw[azulprof,dashed,thick] (\x,-1.4) -- (\x,1.4);
   \node[azulprof,above] at (\x,1.4) {$\v$};
   \fill[laranja] (\x,0) circle (0.07);
   \node[below,font=\footnotesize] at (\x,-0.15) {\lab};}
 % linhas de forca (horizontais com seta)
 \foreach \y in {-1,0,1}
   \draw[Vcol,->,thick] (-0.6,\y) -- (5.0,\y);
 \node[Vcol] at (5.2,0.7) {$\vec{E}$};
\end{tikzpicture}}{Linhas de força (cheias) e equipotenciais (tracejadas).}

Sobre o trabalho da força elétrica nesse deslocamento, é correto afirmar que:
\begin{alternativas}
\item é nulo, pois A e C estão na mesma linha de força.
\item é positivo, pois a carga se move no sentido do campo.
\item é negativo, pois a carga se move contra o campo.
\item depende do caminho seguido entre A e C.
\item é máximo se a carga for levada de A até B.
\end{alternativas}
\resposta{(b)}
\resolucao{A carga positiva se desloca de A (maior potencial) para C (menor
potencial), ou seja, \emph{a favor} do campo elétrico. O trabalho da força
elétrica é $W = q(V_A - V_C) > 0$. O trabalho \emph{não} depende do caminho (força
conservativa), o que elimina (d).}
\end{questao}

\begin{questao}[2]{}
A figura representa as linhas equipotenciais e as linhas de força do campo gerado
por uma carga puntiforme $Q$ fixa no vácuo. As equipotenciais são círculos
concêntricos.

\circuitoq{%
\begin{tikzpicture}[scale=0.9,line width=0.8pt,>=Stealth]
 \shade[ball color=Vcol!30] (0,0) circle (0.28);
 \node at (0,0) {\footnotesize$Q$};
 \foreach \r in {1,1.7,2.4}
   \draw[azulprof,dashed] (0,0) circle (\r);
 \foreach \a in {0,45,...,315}
   \draw[Vcol,->,thick] (\a:0.35) -- (\a:2.7);
 \fill[laranja] (0:1) circle (0.06) node[above right,font=\footnotesize]{A};
 \fill[laranja] (0:2.4) circle (0.06) node[above right,font=\footnotesize]{C};
\end{tikzpicture}}{Campo e equipotenciais de uma carga puntiforme positiva.}

Como as linhas de força \emph{apontam para fora} de $Q$, pode-se concluir que:
\begin{alternativas}
\item $Q$ é negativa e o potencial aumenta com a distância.
\item $Q$ é positiva e o potencial diminui com a distância.
\item $Q$ é positiva e o potencial aumenta com a distância.
\item o potencial é o mesmo em A e C.
\item o campo é uniforme.
\end{alternativas}
\resposta{(b)}
\resolucao{Linhas de força saindo indicam carga \textbf{positiva}. O potencial de
uma carga puntiforme, $V = \ke Q/d$, \emph{diminui} à medida que a distância
cresce: $V_A > V_C$. As equipotenciais, mais afastadas, têm menor potencial. O
campo não é uniforme --- é radial e enfraquece com $1/d^2$.}
\end{questao}

\begin{questao}[2]{}
A figura mostra a configuração das linhas de força do campo gerado por duas
partículas carregadas A e B. As linhas \emph{saem} de A e \emph{entram} em B.

\circuitoq{%
\begin{tikzpicture}[scale=0.85,line width=0.7pt,>=Stealth]
 \shade[ball color=Vcol!25] (-2,0) circle (0.3); \node at (-2,0){\footnotesize A};
 \shade[ball color=Icol!25] ( 2,0) circle (0.3); \node at (2,0){\footnotesize B};
 \foreach \a in {30,60,90,120,150,-30,-60,-90,-120,-150}
   \draw[Vcol,->] (-2,0) .. controls (\a:1.3) and ({180-\a}:1.3) .. (2,0);
 \draw[Vcol,->] (-2,0)--(-3.3,0);
 \draw[Vcol,->] (3.3,0)--(2,0);
\end{tikzpicture}}{Linhas de força de duas cargas A e B.}

Os sinais das cargas A e B são, respectivamente:
\begin{alternativas}[2]
\item positivo e positivo.
\item positivo e negativo.
\item negativo e positivo.
\item negativo e negativo.
\item ambos nulos.
\end{alternativas}
\resposta{(b)}
\resolucao{As linhas de força \emph{nascem} nas cargas positivas e \emph{morrem}
nas negativas. Como saem de A, esta é \textbf{positiva}; como entram em B, esta é
\textbf{negativa}. A configuração é a de um dipolo elétrico.}
\end{questao}

\begin{questao}[2]{}
Uma carga elétrica cria, em um ponto P situado a $\SI{20}{\centi\meter}$ dela, um
campo de intensidade $\SI{900}{\newton\per\coulomb}$. O módulo do potencial
elétrico em P é:
\begin{alternativas}[3]
\item \SI{45}{\volt}
\item \SI{90}{\volt}
\item \SI{180}{\volt}
\item \SI{450}{\volt}
\item \SI{900}{\volt}
\end{alternativas}
\resposta{(c)}
\resolucao{Para uma carga puntiforme, $E = \dfrac{\ke Q}{d^2}$ e
$V = \dfrac{\ke Q}{d}$. Dividindo, $\dfrac{V}{E} = d$, logo
\[
V = E\cdot d = 900\cdot0{,}20 = \SI{180}{\volt}.
\]
Essa relação $V = E\,d$ vale para carga puntiforme; não confundir com o campo
uniforme, onde também $U = E\,d$, mas por outro motivo.}
\end{questao}

\begin{questao}[2]{}
Próximo ao solo, em um dia limpo, a atmosfera apresenta um campo elétrico
aproximadamente uniforme, vertical, apontando para baixo, de intensidade
$\SI{120}{\volt\per\meter}$. Considere dois pontos M (mais alto) e N (mais baixo),
separados verticalmente por $\SI{10}{\meter}$.

\circuitoq{%
\begin{tikzpicture}[scale=0.9,line width=0.8pt,>=Stealth]
 \fill[cinza!25] (-2.5,-1.6) rectangle (2.5,-1.2);
 \node[cinza,font=\footnotesize] at (0,-1.4) {solo};
 \foreach \x in {-2,-1,0,1,2}
   \draw[Vcol,->,thick] (\x,1.4) -- (\x,-1.1);
 \node[Vcol] at (2.5,0.4) {$\vec{E}$};
 \fill[laranja] (-1,1.1) circle (0.08) node[left,font=\footnotesize]{M};
 \fill[laranja] (-1,-0.9) circle (0.08) node[left,font=\footnotesize]{N};
 \draw[<->,cinza] (0.6,1.1)--(0.6,-0.9) node[midway,right,font=\footnotesize]{$\SI{10}{\meter}$};
\end{tikzpicture}}{Campo elétrico atmosférico próximo ao solo.}

A diferença de potencial entre M e N vale:
\begin{alternativas}[2]
\item \SI{12}{\volt}
\item \SI{120}{\volt}
\item \SI{1200}{\volt}
\item \SI{12000}{\volt}
\item zero
\end{alternativas}
\resposta{(c)}
\resolucao{Em campo uniforme, $U = E\,d = 120\cdot10 = \SI{1200}{\volt}$. Como o
campo aponta para baixo (de M para N), M está a maior potencial:
$V_M - V_N = \SI{1200}{\volt}$.}
\end{questao}

\blocoaprofundamento

\begin{questao}[3]{}
Duas cargas $Q_1 = \SI{+2}{\micro\coulomb}$ e $Q_2 = \SI{-2}{\micro\coulomb}$
estão fixas, separadas por $\SI{20}{\centi\meter}$. Determine o potencial elétrico
no ponto médio M do segmento que as une.
\dados{$\ke = \SI{9e9}{\newton\meter\squared\per\coulomb\squared}$.}
\resposta{$V_M = 0$}
\resolucao{O potencial é uma grandeza \emph{escalar}: soma-se com sinal. No ponto
médio, cada carga está a $d = \SI{10}{\centi\meter}$:
\[
V_M = \ke\frac{Q_1}{d}+\ke\frac{Q_2}{d}
= \ke\frac{(+2)+(-2)}{d}\,\si{\micro\coulomb}=0.
\]
O potencial é nulo, mas o \emph{campo} em M \emph{não} é zero --- os vetores campo
das duas cargas apontam no mesmo sentido e se somam. Este é o ponto central: campo
nulo e potencial nulo são condições independentes.}
\end{questao}

\begin{questao}[3]{}
Uma carga puntiforme $Q = \SI{+5}{\micro\coulomb}$ está fixa no vácuo. Uma segunda
carga $q = \SI{+2}{\micro\coulomb}$ é trazida do infinito até um ponto a
$\SI{30}{\centi\meter}$ de $Q$.
\begin{itensq}
\item Qual é o potencial criado por $Q$ nesse ponto?
\item Qual é o trabalho realizado pela força externa para trazer $q$ (supondo-a
      trazida lentamente)?
\end{itensq}
\dados{$\ke = \SI{9e9}{\newton\meter\squared\per\coulomb\squared}$;
$V_\infty = 0$.}
\resposta{(a) $V = \SI{1.5e5}{\volt}$; (b) $W = \SI{0.3}{\joule}$}
\resolucao{(a) $V = \ke\dfrac{Q}{d}
= 9\cdot10^9\cdot\dfrac{\pot{5}{-6}}{0{,}3}=\pot{1.5}{5}\,\si{\volt}$.\\
(b) O trabalho da força externa iguala a energia potencial adquirida (partindo do
infinito, onde $V=0$):
\[
W = q\,(V - V_\infty) = \pot{2}{-6}\cdot\pot{1.5}{5} = \SI{0.3}{\joule}.
\]
Como ambas as cargas são positivas, elas se repelem, e o trabalho externo é
\emph{positivo} (precisamos empurrar $q$ contra a repulsão). Se as cargas tivessem
sinais opostos, o trabalho externo seria negativo.}
\end{questao}

\begin{questao}[3]{}
Um elétron ($q = \SI{-1.6e-19}{\coulomb}$, $m = \SI{9.1e-31}{\kilogram}$) parte do
repouso e é acelerado por uma ddp de $\SI{200}{\volt}$. Determine a energia
cinética adquirida e a velocidade final do elétron.
\resposta{$E_c = \SI{3.2e-17}{\joule}$; $v \approx \SI{8.4e6}{\meter\per\second}$}
\resolucao{Toda a energia potencial elétrica se converte em energia cinética:
\[
E_c = |q|\,U = \pot{1.6}{-19}\cdot200 = \pot{3.2}{-17}\,\si{\joule}.
\]
Da definição de energia cinética, $E_c = \tfrac12 m v^2$:
\[
v=\sqrt{\frac{2E_c}{m}}=\sqrt{\frac{2\cdot\pot{3.2}{-17}}{\pot{9.1}{-31}}}
\approx\SI{8.4e6}{\meter\per\second}.
\]
É o princípio de funcionamento dos antigos tubos de raios catódicos e dos
aceleradores de partículas: usar uma ddp para acelerar cargas.}
\end{questao}

\begin{questao}[3]{}
O diagrama representa o potencial elétrico $V$ criado por uma carga puntiforme em
função da distância $d$ até ela.

\circuitoqq{%
\begin{tikzpicture}
 \begin{axis}[width=9cm,height=5cm,xlabel={$d$ (m)},ylabel={$V$ (V)},
   xmin=0,xmax=6,ymin=0,ymax=90,xtick={0,1,2,3,4},ytick={0,20,40,80},
   grid=both,axis lines=left,domain=0.9:6,samples=100]
   \addplot[Vcol,very thick] {80/x};
   \addplot[only marks,mark=*,Vcol] coordinates {(1,80)(2,40)(4,20)};
   \node[Vcol,anchor=south west,font=\footnotesize] at (axis cs:2,40){$(2,\,40)$};
 \end{axis}
\end{tikzpicture}}{Potencial de uma carga puntiforme em função da distância.}

Sabendo que, a $\SI{2}{\meter}$, o potencial vale $\SI{40}{\volt}$, o potencial a
$\SI{4}{\meter}$ e o campo elétrico a $\SI{2}{\meter}$ valem, respectivamente:
\begin{alternativas}[2]
\item \SI{20}{\volt} e \SI{20}{\volt\per\meter}
\item \SI{20}{\volt} e \SI{40}{\volt\per\meter}
\item \SI{10}{\volt} e \SI{10}{\volt\per\meter}
\item \SI{40}{\volt} e \SI{40}{\volt\per\meter}
\item \SI{10}{\volt} e \SI{20}{\volt\per\meter}
\end{alternativas}
\resposta{(a)}
\resolucao{O potencial varia com $1/d$: dobrando a distância (de \SI{2}{} para
\SI{4}{\meter}), ele cai à metade, $V(4) = \SI{20}{\volt}$.\\
Para o campo, usamos $E = \dfrac{V}{d}$ (válido para carga puntiforme):
$E(2) = \dfrac{40}{2} = \SI{20}{\volt\per\meter}$.}
\end{questao}

\begin{questao}[3]{}
Na figura, uma partícula de carga $Q > 0$ está fixa em P. As superfícies A, B e C
são equipotenciais, e um ponto D situa-se sobre uma linha de força entre B e C.

\circuitoq{%
\begin{tikzpicture}[scale=0.9,line width=0.8pt,>=Stealth]
 \shade[ball color=Vcol!30] (0,0) circle (0.25); \node at (0,0){\footnotesize P};
 \foreach \r/\lab in {1.1/A,1.8/B,2.5/C}{
   \draw[azulprof,dashed] (0,0) circle (\r);
   \node[azulprof,font=\footnotesize] at (\r,0.28) {\lab};}
 \draw[Vcol,->,thick] (20:0.3)--(20:2.7);
 \fill[laranja] (20:2.15) circle (0.07) node[right,font=\footnotesize]{D};
\end{tikzpicture}}{Equipotenciais de uma carga puntiforme positiva.}

Sobre os potenciais, é correto afirmar que:
\begin{alternativas}
\item $V_A < V_B < V_C$.
\item $V_A > V_B > V_C$.
\item $V_A = V_B = V_C$.
\item $V_C$ é o maior, por estar mais longe.
\item nada se pode afirmar sem o valor de $Q$.
\end{alternativas}
\resposta{(b)}
\resolucao{Para carga positiva, $V = \ke Q/d$ \emph{decresce} com a distância. Como
A é a superfície mais interna (menor $d$) e C a mais externa (maior $d$):
$V_A > V_B > V_C$. O sinal de $Q$ determina a tendência; o \emph{valor} de $Q$ não
é necessário para ordenar os potenciais.}
\end{questao}

\begin{questao}[3]{}
Uma esfera condutora de raio $R = \SI{10}{\centi\meter}$ está carregada com
$Q = \SI{6}{\micro\coulomb}$, isolada no vácuo. Os gráficos mostram o campo
elétrico e o potencial em função da distância $d$ ao centro.

\circuitoqq{%
\begin{tikzpicture}
 \begin{axis}[width=6.2cm,height=4.2cm,xlabel={$d$ (m)},ylabel={$E$ (kN/C)},
   xmin=0,xmax=0.45,ymin=0,ymax=6200,xtick={0,0.1,0.2,0.3,0.4},
   ytick={0,2000,4000,5400},yticklabels={0,2000,4000,5400},
   grid=both,axis lines=left,domain=0.1:0.45,samples=80,
   scaled y ticks=false,y tick label style={/pgf/number format/fixed}]
   \addplot[Vcol,very thick] {9e9*6e-6/(x^2)/1000};
   \addplot[Vcol,very thick] coordinates {(0,0)(0.1,0)};
   \addplot[Vcol,dashed] coordinates {(0.1,0)(0.1,5400)};
 \end{axis}
\end{tikzpicture}\hspace{4mm}
\begin{tikzpicture}
 \begin{axis}[width=6.2cm,height=4.2cm,xlabel={$d$ (m)},ylabel={$V$ (kV)},
   xmin=0,xmax=0.45,ymin=0,ymax=620,xtick={0,0.1,0.2,0.3,0.4},
   ytick={0,180,540},grid=both,axis lines=left,domain=0.1:0.45,samples=80]
   \addplot[Icol,very thick] {9e9*6e-6/x/1000};
   \addplot[Icol,very thick] coordinates {(0,540)(0.1,540)};
 \end{axis}
\end{tikzpicture}}{Campo e potencial de uma esfera condutora carregada.}

\begin{itensq}
\item Calcule $E$ e $V$ na superfície da esfera.
\item Determine $E$ e $V$ a $\SI{30}{\centi\meter}$ do centro.
\item Explique por que, \emph{dentro} da esfera, $E = 0$ mas $V \neq 0$.
\end{itensq}
\dados{$\ke = \SI{9e9}{\newton\meter\squared\per\coulomb\squared}$.}
\resposta{(a) $E=\SI{5.4e6}{\newton\per\coulomb}$, $V=\SI{540}{\kilo\volt}$;
(b) $E=\SI{6e5}{\newton\per\coulomb}$, $V=\SI{180}{\kilo\volt}$; (c) ver comentário}
\resolucao{(a) Fora da esfera (inclusive na superfície), ela se comporta como uma
carga puntiforme concentrada no centro:
\[
E(R)=\frac{\ke Q}{R^{2}}=\frac{9\cdot10^9\cdot\pot{6}{-6}}{(0{,}1)^2}
=\pot{5.4}{6}\,\si{\newton\per\coulomb},
\]
\[
V(R)=\frac{\ke Q}{R}=\frac{9\cdot10^9\cdot\pot{6}{-6}}{0{,}1}
=\pot{5.4}{5}\,\si{\volt}=\SI{540}{\kilo\volt}.
\]
(b) A $\SI{30}{\centi\meter}$ (três vezes o raio): $E$ cai com $1/d^2$ (fator 9) e
$V$ com $1/d$ (fator 3):
\[
E=\frac{5{,}4\cdot10^6}{9}=\pot{6}{5}\,\si{\newton\per\coulomb},
\qquad
V=\frac{540}{3}=\SI{180}{\kilo\volt}.
\]
(c) \textbf{A distinção conceitual central.} Em um condutor em equilíbrio
eletrostático, toda a carga se distribui na \emph{superfície} e o campo interno é
nulo --- se não fosse, as cargas livres se moveriam e não haveria equilíbrio.\\
Mas o potencial \emph{não} é nulo: ele é a \emph{integral} do campo desde o
infinito até o ponto. Ao entrar na esfera, o campo passa a ser zero, então o
potencial \textbf{para de variar} --- ele permanece constante e igual ao valor da
superfície:
\[
V_{\text{interno}} = V(R) = \SI{540}{\kilo\volt}.
\]
Isso é coerente com $E = -\dfrac{\mathrm{d}V}{\mathrm{d}d}$: campo nulo significa
potencial \emph{constante}, não necessariamente nulo. Todo o volume do condutor é
uma única região equipotencial.\\
\textbf{Aplicação:} é o princípio da \textbf{gaiola de Faraday} --- dentro de um
condutor fechado o campo é nulo, protegendo o conteúdo, ainda que o conjunto esteja
a um potencial elevadíssimo. É por isso que se está seguro dentro de um carro
atingido por um raio.}
\end{questao}

\begin{questao}[3]{}
Três cargas puntiformes iguais, $q = \SI{1}{\micro\coulomb}$, são trazidas
\emph{uma a uma} desde o infinito até ocuparem os vértices de um triângulo
equilátero de lado $L = \SI{10}{\centi\meter}$.
\begin{itensq}
\item Calcule o trabalho externo necessário para trazer cada carga.
\item Determine a energia potencial total da configuração.
\item Se as cargas forem liberadas simultaneamente, qual é a energia cinética total
      quando estiverem infinitamente afastadas?
\end{itensq}
\dados{$\ke = \SI{9e9}{\newton\meter\squared\per\coulomb\squared}$.}
\resposta{(a) $0$, \SI{0.09}{\joule}, \SI{0.18}{\joule};
(b) $U = \SI{0.27}{\joule}$; (c) \SI{0.27}{\joule}}
\resolucao{\textbf{(a) Trabalho carga a carga.}
\begin{itemize}[leftmargin=1.6em,itemsep=1pt,topsep=2pt]
\item \textbf{1\textsuperscript{a} carga:} não há campo algum ainda, logo
      $W_1 = 0$.
\item \textbf{2\textsuperscript{a} carga:} vem contra o campo da primeira, a uma
      distância final $L$:
      \[
      W_2=\ke\frac{q^2}{L}=9\cdot10^9\frac{(\pot{1}{-6})^2}{0{,}1}=\SI{0.09}{\joule}.
      \]
\item \textbf{3\textsuperscript{a} carga:} sofre a ação das \emph{duas} anteriores,
      ambas à distância $L$:
      \[
      W_3=2\,\ke\frac{q^2}{L}=\SI{0.18}{\joule}.
      \]
\end{itemize}
\textbf{(b) Energia total da configuração.} É a soma dos trabalhos --- ou,
equivalentemente, a soma sobre os \emph{três pares} distintos:
\[
U = W_1+W_2+W_3 = 0+0{,}09+0{,}18 = \SI{0.27}{\joule}
= 3\,\ke\frac{q^2}{L}. \checkmark
\]
Note que a resposta \emph{não depende} da ordem em que as cargas foram trazidas
--- consequência de a força elétrica ser conservativa.

\textbf{(c)} Ao serem liberadas, toda a energia potencial armazenada se converte
em cinética (conservação de energia, sem dissipação):
\[
E_{c,\text{total}} = U = \SI{0.27}{\joule}.
\]
\textbf{Cuidado com o erro clássico:} a energia total \emph{não} é
$3\times W_3$ nem a soma de "cada carga vezes o potencial das outras" sem dividir
por 2 --- isso contaria cada par duas vezes. A fórmula geral é
$U=\tfrac12\sum_i q_i V_i$, onde $V_i$ é o potencial criado \emph{pelas demais}
na posição de $i$.}
\end{questao}

\begin{questao}[3]{}
Duas esferas condutoras, de raios $R$ e $3R$, estão distantes entre si e são
ligadas por um fio condutor longo e fino. A carga total do sistema é $Q$.
\begin{itensq}
\item Determine a carga final de cada esfera.
\item Compare os campos elétricos nas superfícies das duas esferas.
\item Explique o "poder das pontas" a partir desse resultado.
\end{itensq}
\resposta{(a) $Q_1=Q/4$, $Q_2=3Q/4$; (b) $E_1/E_2=3$; (c) ver comentário}
\resolucao{(a) Ligadas por um condutor, as esferas atingem o \emph{mesmo
potencial}:
\[
\frac{\ke Q_1}{R}=\frac{\ke Q_2}{3R}
\;\Longrightarrow\;
Q_2 = 3Q_1.
\]
Com $Q_1+Q_2=Q$: $Q_1=\dfrac{Q}{4}$ e $Q_2=\dfrac{3Q}{4}$.\\
(b) Campos nas superfícies:
\[
E_1=\frac{\ke Q_1}{R^{2}}=\frac{\ke Q}{4R^{2}},
\qquad
E_2=\frac{\ke Q_2}{(3R)^{2}}=\frac{3\ke Q}{4\cdot9R^{2}}=\frac{\ke Q}{12R^{2}}.
\]
\[
\frac{E_1}{E_2}=3.
\]
A esfera \textbf{menor}, embora tenha \emph{menos} carga, apresenta campo
\textbf{três vezes mais intenso} em sua superfície.\\
\textbf{Relação geral:} como $V=\ke Q/R$ e $E=\ke Q/R^{2}$, tem-se
$E = V/R$. Com $V$ comum às duas, \emph{o campo é inversamente proporcional ao
raio}.\\
(c) \textbf{Poder das pontas.} Uma "ponta" é uma região de raio de curvatura
muito pequeno. Sendo $E \propto 1/R$, o campo ali fica extremamente intenso ---
suficiente para ionizar o ar e provocar descarga, mesmo com o objeto a um
potencial moderado.\\
\textbf{Consequências práticas:}
\begin{itemize}[leftmargin=1.6em,itemsep=1pt,topsep=2pt]
\item \textbf{Para-raios} são pontiagudos \emph{de propósito}: concentram o campo
      e induzem a descarga em local controlado.
\item Equipamentos de alta tensão têm cantos \emph{arredondados} e cúpulas
      polidas, justamente para \emph{evitar} campos intensos e o efeito corona.
\end{itemize}}
\end{questao}

\blocodesafio

\begin{questao}[4]{}
Quatro cargas iguais $q=\SI{1}{\micro\coulomb}$ ocupam os vértices de um quadrado
de lado $a=\SI{10}{\centi\meter}$.
\begin{itensq}
\item Calcule a energia potencial elétrica total da configuração.
\item Quanto trabalho externo seria necessário para trazer uma quinta carga $q$
      desde o infinito até o centro do quadrado?
\item Se todas fossem liberadas, qual seria a energia cinética total no infinito?
\end{itensq}
\dados{$\ke=\SI{9e9}{}$; $\sqrt2\approx1{,}41$.}
\resposta{(a) $U\approx\SI{0.49}{\joule}$; (b) $W\approx\SI{0.51}{\joule}$;
(c) $\approx\SI{0.49}{\joule}$}
\resolucao{(a) Há $\binom{4}{2}=6$ pares: \textbf{4 lados} (distância $a$) e
\textbf{2 diagonais} (distância $a\sqrt2$):
\[
U=4\cdot\frac{\ke q^{2}}{a}+2\cdot\frac{\ke q^{2}}{a\sqrt2}
=\frac{\ke q^{2}}{a}\left(4+\frac{2}{\sqrt2}\right)
=\frac{\ke q^{2}}{a}\left(4+\sqrt2\right).
\]
Com $\dfrac{\ke q^{2}}{a}=\dfrac{9\cdot10^{9}(\pot{1}{-6})^{2}}{0{,}1}
=\SI{0.09}{\joule}$:
\[
U=0{,}09\,(4+1{,}41)\approx\SI{0.49}{\joule}.
\]
(b) A quinta carga, no centro, dista $\dfrac{a\sqrt2}{2}=\SI{0.0707}{\meter}$ de
\emph{cada} uma das quatro. O potencial no centro é
\[
V_c=4\cdot\frac{\ke q}{a\sqrt2/2}
=4\cdot\frac{9\cdot10^{9}\cdot\pot{1}{-6}}{0{,}0707}
\approx\pot{5.09}{5}\,\si{\volt},
\]
e o trabalho externo é
\[
W=q\,V_c=\pot{1}{-6}\cdot\pot{5.09}{5}\approx\SI{0.51}{\joule}.
\]
(c) Sem a quinta carga, ao liberar as quatro, toda a energia potencial da
configuração converte-se em cinética:
\[
E_{c}=U\approx\SI{0.49}{\joule}.
\]
\textbf{Atenção ao método.} Em (a) somam-se os \emph{pares} (cada par uma vez
só); em (b) usa-se $W=qV$ porque as quatro cargas já estavam posicionadas e
imóveis. Confundir os dois procedimentos --- contando pares duas vezes --- é o
erro mais comum neste tipo de problema.}
\end{questao}

\begin{questao}[4]{}
Uma esfera condutora oca de raio interno $a$ e raio externo $b$ envolve, no
centro, uma carga puntiforme $+Q$. A casca está inicialmente \emph{neutra} e
isolada.
\begin{itensq}
\item Determine as cargas induzidas nas superfícies interna e externa da casca.
\item Esboce o comportamento de $E(r)$ nas quatro regiões ($r<a$, $a<r<b$,
      $r>b$).
\item Se a casca for \textbf{aterrada}, o que muda?
\end{itensq}
\resposta{(a) $-Q$ na interna, $+Q$ na externa; (b) e (c) ver comentário}
\resolucao{(a) Dentro do \emph{material} do condutor ($a<r<b$) o campo deve ser
\textbf{nulo} em equilíbrio. Aplicando a lei de Gauss a uma superfície esférica
nessa região, a carga interna total deve ser zero:
\[
Q + q_{\text{int}} = 0 \;\Longrightarrow\; q_{\text{int}}=-Q.
\]
Como a casca é neutra no total, a superfície externa fica com
$q_{\text{ext}}=+Q$.\\
(b) \textbf{Comportamento do campo:}
\begin{center}\small\sffamily
\begin{tabular}{@{}l l l@{}}
\toprule
\textbf{Região} & \textbf{Carga envolvida} & \textbf{Campo}\\
\midrule
$r<a$      & $+Q$        & $E=\ke Q/r^{2}$ (radial, para fora)\\
$a<r<b$    & $Q-Q=0$     & $E=0$ (dentro do condutor)\\
$r>b$      & $Q-Q+Q=+Q$  & $E=\ke Q/r^{2}$\\
\bottomrule
\end{tabular}
\end{center}
Fora da casca, o campo é \emph{idêntico} ao que existiria sem ela --- a casca
neutra não blinda o exterior.\\
(c) \textbf{Aterrando a casca}, a superfície externa perde sua carga $+Q$ para a
terra (os elétrons sobem para neutralizá-la). Resultado:
\[
q_{\text{int}}=-Q,\qquad q_{\text{ext}}=0
\;\Longrightarrow\;
E = 0 \ \text{ para } r>b.
\]
\textbf{Agora sim há blindagem completa}: nenhum campo escapa para fora.\\
\textbf{Aplicação --- blindagem eletrostática.} É por isso que cabos coaxiais e
gabinetes de equipamentos sensíveis têm sua malha metálica \textbf{aterrada}:
uma blindagem flutuante (não aterrada) não impede o campo de vazar. O aterramento
é o que transforma o condutor em uma verdadeira gaiola de Faraday.}
\end{questao}

% BEGIN BANCO COMPLEMENTAR Potencial elétrico e superfícies equipotenciais
% =====================================================================
%  Banco complementar --- Potencial Eletrico e Superficies
%  Equipotenciais  (REVISADO)
%  Substitui a versao gerada por template: as 6 questoes N1 eram o
%  mesmo enunciado ("determine V a r de uma carga Q") com numeros
%  trocados, as 5 N2 eram "carga q atravessa uma ddp" repetido, e as
%  4 N3 eram "no ponto P as distancias as cargas q1 e q2 sao..."
%  igualmente clonado. As 6 N4 eram frases conceituais de uma linha.
%  O cap01c ja cobre: equipotenciais (MC), campo uniforme entre A e B,
%  tres equipotenciais, dipolo com linhas, linhas de forca com
%  equipotenciais, carga puntiforme, duas particulas A e B, campo a
%  20 cm, campo atmosferico, cargas +-2 uC, trabalho, eletron
%  acelerado, grafico V x d, superficies A/B/C, esfera condutora de
%  R=10 cm, tres cargas trazidas do infinito, duas esferas ligadas por
%  fio, quadrado de quatro cargas e esfera oca com carga central.
%  Distribuicao mantida: 6 N1 + 5 N2 + 4 N3 + 6 N4 = 21
% =====================================================================

\subsubsection*{Banco complementar --- Potencial elétrico e superfícies equipotenciais}

\begin{atencao}[Organização do banco]
O potencial é \textbf{escalar}: contribuições de várias cargas somam-se
algebricamente, com sinal, sem decomposição vetorial --- o que torna $V$ muito
mais fácil de calcular que $\vec E$. Em compensação, $V$ carrega menos
informação: pontos de $V=0$ podem ter campo intenso, e é o \emph{gradiente},
não o valor, que determina a força. O N4 explora essa distinção, a
independência do caminho e a estabilidade.
\end{atencao}

\blocofixacao

\begin{questao}[1]{}
Determine o potencial a $\SI{30}{\centi\meter}$ de uma carga
$Q=+\SI{2}{\micro\coulomb}$, adotando $V(\infty)=0$.
\dados{$\ke=\SI{9e9}{\newton\meter\squared\per\coulomb\squared}$}
\resposta{$V=\SI{60}{\kilo\volt}$}
\resolucao{
\[
V=\frac{\ke Q}{r}=\frac{\pot{9}{9}\times\pot{2}{-6}}{0{,}30}
=\frac{\pot{1.8}{4}}{0{,}30}=\pot{6}{4}\,\si{\volt}=\SI{60}{\kilo\volt}.
\]
\textbf{Atenção à potência de $r$:} o potencial cai com $1/r$, enquanto o campo
cai com $1/r^{2}$. Confundir as duas é o erro mais comum do tópico --- e o
resultado sai com a dimensão errada.}
\end{questao}

\begin{questao}[1]{}
Determine o potencial a $\SI{60}{\centi\meter}$ de uma carga
$Q=-\SI{4}{\micro\coulomb}$.
\resposta{$V=\SI{-60}{\kilo\volt}$}
\resolucao{
\[
V=\frac{\pot{9}{9}\times(-\pot{4}{-6})}{0{,}60}=-\pot{6}{4}\,\si{\volt}
=\SI{-60}{\kilo\volt}.
\]
\textbf{O sinal é do potencial, não do módulo.} Diferentemente do campo --- cuja
\emph{orientação} é dada por um vetor --- o potencial é um escalar com sinal
próprio. Cargas negativas produzem potencial negativo em toda parte.}
\end{questao}

\begin{questao}[1]{}
A que distância de $Q=+\SI{3}{\micro\coulomb}$ o potencial vale
$\SI{45}{\kilo\volt}$?
\resposta{$r=\SI{0.6}{\meter}$}
\resolucao{
\[
r=\frac{\ke Q}{V}=\frac{\pot{9}{9}\times\pot{3}{-6}}{\pot{4.5}{4}}
=\frac{\pot{2.7}{4}}{\pot{4.5}{4}}=\SI{0.6}{\meter}.
\]
\textbf{Superfície equipotencial:} \emph{todos} os pontos a
$\SI{0.6}{\meter}$ da carga estão a $\SI{45}{\kilo\volt}$. Em torno de uma
carga puntiforme, as equipotenciais são esferas concêntricas.}
\end{questao}

\begin{questao}[1]{}
Determine a carga que produz $\SI{-18}{\kilo\volt}$ a
$\SI{0.5}{\meter}$ de distância.
\resposta{$Q=-\SI{1}{\micro\coulomb}$}
\resolucao{
\[
Q=\frac{Vr}{\ke}=\frac{(-\pot{1.8}{4})(0{,}5)}{\pot{9}{9}}
=-\pot{1}{-6}\,\si{\coulomb}=-\SI{1}{\micro\coulomb}.
\]
O sinal negativo do potencial revelou diretamente o sinal da carga --- vantagem
prática de trabalhar com uma grandeza escalar.}
\end{questao}

\begin{questao}[1]{}
Sobre o potencial criado por várias cargas em um ponto, é correto afirmar:
\begin{alternativas}[2]
\item soma-se vetorialmente
\item soma-se algebricamente, com sinal
\item soma-se em módulo
\item depende do caminho até o ponto
\end{alternativas}
\resposta{(b)}
\resolucao{O potencial é \textbf{escalar}: basta somar as contribuições
$\ke q_k/r_k$ com seus sinais, sem ângulos nem componentes.\\
\textbf{Vantagem prática enorme:} calcular $V$ de dez cargas exige dez divisões
e uma soma; calcular $\vec E$ exigiria dez vetores decompostos em componentes.\\
A alternativa (d) descreve a propriedade oposta --- justamente por o campo
eletrostático ser \emph{conservativo}, o potencial de um ponto independe do
caminho, e é isso que permite defini-lo.}
\end{questao}

\begin{questao}[1]{}
Uma carga $q=+\SI{3}{\micro\coulomb}$ é colocada em um ponto onde
$V=\SI{200}{\volt}$. Determine sua energia potencial elétrica.
\resposta{$U=\SI{0.6}{\milli\joule}$}
\resolucao{
\[
U=qV=\pot{3}{-6}\times200=\pot{6}{-4}\,\si{\joule}=\SI{0.6}{\milli\joule}.
\]
\textbf{Distinção essencial:} $V$ é uma propriedade \emph{do ponto} --- existe
mesmo sem carga alguma ali. Já $U$ é uma propriedade \emph{do sistema}
carga-campo, e só faz sentido quando a carga está presente.\\
A relação $U=qV$ é o análogo elétrico de $E_p=mgh$: potencial faz o papel de
$gh$, e carga faz o papel de massa.}
\end{questao}

\blocoaplicacao

\begin{questao}[2]{}
Uma carga $q=+\SI{2}{\micro\coulomb}$ desloca-se de um ponto a
$\SI{50}{\volt}$ para outro a $\SI{20}{\volt}$.
\begin{itensq}
\item Determine o trabalho realizado pela força elétrica.
\item Determine a variação da energia potencial.
\end{itensq}
\resposta{$W=\SI{60}{\micro\joule}$; $\Delta U=\SI{-60}{\micro\joule}$}
\resolucao{(a)
\[
W=q\left(V_A-V_B\right)=\pot{2}{-6}(50-20)=\pot{6}{-5}\,\si{\joule}
=\SI{60}{\micro\joule}.
\]
(b) $\Delta U=-W=\SI{-60}{\micro\joule}$.\\
\textbf{Coerência dos sinais:} a carga é positiva e foi de maior para menor
potencial --- movimento espontâneo. A força elétrica realiza trabalho
\emph{positivo} e a energia potencial \emph{diminui}, convertendo-se em energia
cinética.\\
\textbf{Regra:} cargas positivas "descem" o potencial espontaneamente; cargas
negativas "sobem".}
\end{questao}

\begin{questao}[2]{}
Duas cargas estão fixas: $q_1=+\SI{3}{\micro\coulomb}$ a
$\SI{30}{\centi\meter}$ do ponto $P$, e $q_2=-\SI{2}{\micro\coulomb}$ a
$\SI{20}{\centi\meter}$ de $P$. Determine o potencial em $P$.
\resposta{$V=0$}
\resolucao{Somando algebricamente:
\[
V=\frac{\pot{9}{9}(\pot{3}{-6})}{0{,}30}
+\frac{\pot{9}{9}(-\pot{2}{-6})}{0{,}20}
=\pot{9}{4}-\pot{9}{4}=\SI{0}{\volt}.
\]
\textbf{Cuidado:} $V=0$ \textbf{não} significa campo nulo nem ausência de força.
As duas cargas produzem campos que \emph{não} se cancelam --- elas têm sinais
opostos, e seus campos apontam para o mesmo lado ao longo da linha que as une.\\
Uma carga colocada em $P$ teria energia potencial nula, mas sofreria força.
Essa distinção é aprofundada no bloco seguinte.}
\end{questao}

\begin{questao}[2]{}
Em um campo uniforme de $E=\SI{500}{\volt\per\meter}$, dois pontos estão
separados por $\SI{4}{\centi\meter}$ ao longo da direção do campo. Determine a
diferença de potencial.
\resposta{$U=\SI{20}{\volt}$}
\resolucao{
\[
U=E\,d=500\times0{,}04=\SI{20}{\volt},
\]
com o ponto \emph{a montante} do campo em potencial maior --- o campo aponta
sempre no sentido de $V$ decrescente.\\
\textbf{Se o deslocamento fosse perpendicular} ao campo, a diferença seria
\emph{nula}: os dois pontos estariam na mesma equipotencial. Só a componente do
deslocamento \emph{ao longo} do campo contribui.\\
\textbf{Unidades:} $\si{\volt\per\meter}$ e $\si{\newton\per\coulomb}$ são a
mesma coisa --- verifique multiplicando por metro e por coulomb.}
\end{questao}

\begin{questao}[2]{}
Um elétron é acelerado através de uma diferença de potencial de
$\SI{500}{\volt}$.
\begin{itensq}
\item Determine a energia adquirida, em joules e em elétron-volts.
\item Determine sua velocidade final, partindo do repouso.
\end{itensq}
\dados{$e=\pot{1.6}{-19}\,\si{\coulomb}$;
$m_e=\pot{9.11}{-31}\,\si{\kilogram}$}
\resposta{$\SI{500}{\electronvolt}=\pot{8}{-17}\,\si{\joule}$;
$v=\pot{1.33}{7}\,\si{\meter\per\second}$}
\resolucao{(a)
\[
W=eU=\pot{1.6}{-19}(500)=\pot{8}{-17}\,\si{\joule}.
\]
Em elétron-volts, a leitura é imediata: $\SI{500}{\electronvolt}$ --- \emph{por
definição}, $\SI{1}{\electronvolt}$ é a energia que a carga elementar adquire
sob $\SI{1}{\volt}$.\\
(b) Toda a energia vira cinética:
\[
v=\sqrt{\frac{2W}{m}}=\sqrt{\frac{2(\pot{8}{-17})}{\pot{9.11}{-31}}}
=\sqrt{\pot{1.756}{14}}=\pot{1.33}{7}\,\si{\meter\per\second}.
\]
\textbf{Verificação de consistência:} isso é $4{,}4\%$ da velocidade da luz ---
ainda no regime não relativístico, onde $\frac12mv^{2}$ é válido. Acima de cerca
de $\SI{50}{\kilo\electronvolt}$, a correção relativística passa a ser
necessária.}
\end{questao}

\begin{questao}[2]{}
Duas cargas, $q_1=+\SI{4}{\micro\coulomb}$ em $x=0$ e
$q_2=-\SI{1}{\micro\coulomb}$ em $x=\SI{0.5}{\meter}$, estão fixas. Determine o
ponto do segmento entre elas onde o potencial é nulo.
\resposta{$x=\SI{0.4}{\meter}$}
\resolucao{Impondo $V=0$ com $x$ medido a partir de $q_1$:
\[
\frac{\ke(4\mu)}{x}+\frac{\ke(-1\mu)}{0{,}5-x}=0
\;\Longrightarrow\;
\frac{4}{x}=\frac{1}{0{,}5-x}.
\]
\[
4(0{,}5-x)=x
\;\Longrightarrow\;
2=5x
\;\Longrightarrow\;
x=\SI{0.4}{\meter}.
\]
\textbf{Contraste importante com o campo:} o ponto de \emph{campo} nulo dessas
mesmas cargas fica \emph{fora} do segmento (a $\SI{1}{\meter}$ de $q_1$),
porque o campo é vetorial e só pode se cancelar onde os dois vetores forem
opostos. Já o potencial é escalar e se anula sempre que as contribuições tiverem
módulos iguais e sinais opostos --- o que ocorre dentro do segmento.\\
\textbf{Existe também um segundo ponto de $V=0$} fora do segmento, do lado da
carga menor, em $x=\SI{0.667}{\meter}$. A superfície completa de $V=0$ é uma
esfera.}
\end{questao}

\blocoaprofundamento

\begin{questao}[3]{}
Duas cargas $+Q$ e $-Q$, com $Q=\SI{2}{\micro\coulomb}$, estão separadas por
$\SI{40}{\centi\meter}$.
\begin{itensq}
\item Determine o potencial no ponto médio.
\item Determine o campo no ponto médio.
\item Explique a aparente contradição.
\end{itensq}
\resposta{$V=0$; $E=\SI{900}{\kilo\volt\per\meter}$}
\resolucao{(a) As duas contribuições têm módulos iguais e sinais opostos:
\[
V=\frac{\ke Q}{0{,}20}-\frac{\ke Q}{0{,}20}=0.
\]
(b) Os \emph{campos}, ao contrário, apontam ambos da carga positiva para a
negativa --- eles se \textbf{somam}:
\[
E=2\times\frac{\ke Q}{(0{,}20)^{2}}
=2\times\frac{\pot{9}{9}(\pot{2}{-6})}{0{,}04}
=\pot{9}{5}\,\si{\volt\per\meter}.
\]
(c) \textbf{Não há contradição.} O potencial não determina o campo --- quem o
determina é o \emph{gradiente} do potencial:
\[
\vec E=-\nabla V.
\]
Uma função pode valer zero em um ponto e ter derivada grande ali. É exatamente o
caso: caminhando ao longo do eixo, $V$ passa de positivo a negativo cruzando o
zero com inclinação acentuada.\\
\textbf{Analogia:} a altitude do nível do mar é zero, mas isso nada diz sobre a
inclinação do terreno naquele ponto. É a inclinação --- e não a altitude --- que
faz uma bola rolar.\\
\textbf{Consequência:} uma carga solta no ponto médio \emph{não} fica parada,
apesar de $V=0$ e $U=0$.}
\end{questao}

\begin{questao}[3]{}
O potencial ao longo de um eixo é dado por
$V(x)=\SI{200}{}-\SI{50}{}\,x$ (com $V$ em volts e $x$ em metros), no trecho
$0\le x\le\SI{4}{\meter}$.
\begin{itensq}
\item Determine o campo elétrico.
\item Determine o trabalho para levar $q=+\SI{2}{\micro\coulomb}$ de
      $x=1$ a $x=\SI{3}{\meter}$.
\item Identifique as superfícies equipotenciais.
\end{itensq}
\resposta{$E=\SI{50}{\volt\per\meter}$ no sentido $+x$;
$W=\SI{200}{\micro\joule}$}
\resolucao{(a)
\[
E_x=-\frac{\mathrm{d}V}{\mathrm{d}x}=-(-50)=\SI{+50}{\volt\per\meter}.
\]
Campo uniforme, apontando no sentido \emph{crescente} de $x$ --- ou seja, no
sentido em que $V$ \emph{decresce}, como sempre.\\
(b) $V(1)=\SI{150}{\volt}$ e $V(3)=\SI{50}{\volt}$:
\[
W=q\left(V_1-V_3\right)=\pot{2}{-6}(150-50)=\SI{200}{\micro\joule}.
\]
\emph{(Conferindo por força: $W=qE\,d=\pot{2}{-6}(50)(2)
=\SI{200}{\micro\joule}$ \checkmark.)}\\
(c) Como $V$ depende \emph{apenas} de $x$, as equipotenciais são
\textbf{planos perpendiculares ao eixo} $x$. O plano $x=\SI{1}{\meter}$ é toda
a superfície com $V=\SI{150}{\volt}$.\\
\textbf{Regra geral:} equipotenciais são sempre \textbf{perpendiculares} às
linhas de campo. Se não fossem, haveria componente de campo ao longo da
superfície, e mover uma carga sobre ela exigiria trabalho --- contradizendo a
definição de equipotencial.}
\end{questao}

\begin{questao}[3]{}
Uma carga $q=+\SI{1}{\micro\coulomb}$ é trazida do infinito até
$\SI{20}{\centi\meter}$ de uma carga fixa $Q=+\SI{5}{\micro\coulomb}$.
\begin{itensq}
\item Determine o trabalho realizado por um agente externo.
\item Determine a energia potencial do par.
\item Interprete o sinal.
\end{itensq}
\resposta{$W_{ext}=U=\SI{0.225}{\joule}$}
\resolucao{(a) e (b) Como a carga parte do repouso no infinito e chega em
repouso, todo o trabalho externo vira energia potencial:
\[
U=\frac{\ke Qq}{r}=\frac{\pot{9}{9}(\pot{5}{-6})(\pot{1}{-6})}{0{,}20}
=\frac{\pot{4.5}{-2}}{0{,}20}=\SI{0.225}{\joule}.
\]
\emph{(Equivalentemente, $U=qV$ com $V=\ke Q/r=\SI{225}{\kilo\volt}$.)}\\
(c) \textbf{Sinal positivo:} as cargas se repelem, então aproximá-las exige
trabalho \emph{contra} a força elétrica. O agente externo gasta energia, que
fica armazenada no sistema.\\
\textbf{Se as cargas tivessem sinais opostos}, $U$ seria negativa: elas se
atraem, e seria preciso \emph{segurá-las} durante a aproximação --- o agente
externo receberia energia.\\
\textbf{Verificação de coerência:} soltando a carga, ela é repelida e ganha
$\SI{0.225}{\joule}$ de energia cinética ao voltar ao infinito. A energia
potencial é integralmente recuperável --- marca de um campo conservativo.}
\end{questao}

\begin{questao}[3]{}
Compare, para duas cargas positivas iguais, o lugar geométrico dos pontos de
potencial nulo e o dos pontos de campo nulo.
\begin{itensq}
\item Determine onde $V=0$.
\item Determine onde $\vec E=\vec 0$.
\item Explique a diferença.
\end{itensq}
\resposta{$V$ nunca se anula (exceto no infinito); $\vec E$ se anula no ponto
médio}
\resolucao{(a) Ambas as cargas são positivas, logo cada uma contribui com
potencial \textbf{positivo} em todo ponto finito. A soma de dois positivos nunca
é zero:
\[
V=\frac{\ke Q}{r_1}+\frac{\ke Q}{r_2}>0\quad\text{sempre}.
\]
$V\to0$ apenas no limite $r\to\infty$.\\
(b) O campo é vetorial e \emph{pode} se cancelar. No ponto médio do segmento,
os dois vetores têm módulos iguais e sentidos opostos:
\[
\vec E=\vec 0\ \text{no ponto médio}.
\]
(c) \textbf{A diferença é a natureza das grandezas.}
\begin{itemize}
\item $V$ é escalar e, para cargas de mesmo sinal, as contribuições não podem se
cancelar --- só se somam.
\item $\vec E$ é vetorial e se cancela sempre que os vetores forem opostos, o que
não impõe restrição alguma sobre os sinais das cargas.
\end{itemize}
\textbf{Contraste com o caso de cargas opostas} ($+Q$ e $-Q$): ali ocorre
exatamente o inverso --- $V=0$ no ponto médio, mas $\vec E\neq\vec 0$.\\
\textbf{Síntese:} $V=0$ e $\vec E=\vec 0$ são condições \emph{independentes}.
Nenhuma implica a outra, em nenhum sentido.}
\end{questao}

\blocodesafio

\begin{questao}[4]{}
\textbf{(Demonstração --- trabalho e independência do caminho.)} Prove que o
trabalho da força elétrica entre dois pontos independe do caminho, e deduza a
relação com a diferença de potencial.
\begin{itensq}
\item Calcule o trabalho para o campo de uma carga puntiforme.
\item Generalize para uma distribuição qualquer.
\item Explique por que isso permite \emph{definir} o potencial.
\end{itensq}
\resposta{$W_{AB}=q(V_A-V_B)$, com
$V=\ke Q/r$}
\resolucao{(a) Para uma carga fixa $Q$ e uma carga de prova $q$, a força é
radial: $\vec F=\dfrac{\ke Qq}{r^{2}}\hat r$. Em um deslocamento
$\mathrm{d}\vec\ell$, apenas a componente radial contribui, pois
$\vec F\cdot\mathrm{d}\vec\ell=F\,\mathrm{d}r$:
\[
W_{AB}=\int_{r_A}^{r_B}\frac{\ke Qq}{r^{2}}\,\mathrm{d}r
=\ke Qq\left[-\frac{1}{r}\right]_{r_A}^{r_B}
=\ke Qq\left(\frac{1}{r_A}-\frac{1}{r_B}\right).
\]
\textbf{O resultado depende apenas de $r_A$ e $r_B$} --- as componentes
transversais do caminho não contribuem, por serem perpendiculares à força.\\
(b) \textbf{Generalização:} para $n$ cargas, a força total é a soma vetorial das
forças individuais, e o trabalho total é a soma dos trabalhos. Como cada parcela
independe do caminho, a soma também independe.\\
Definindo $V=\sum\ke Q_k/r_k$, obtém-se
\[
\boxed{W_{AB}=q\left(V_A-V_B\right)}
\]
(c) \textbf{Por que isso permite definir $V$.} Se o trabalho dependesse do
caminho, não haveria como atribuir um \emph{número} a cada ponto --- o "potencial"
seria ambíguo. A independência do caminho garante que
$V(P)=W_{\infty\to P}/q$ é bem definido.\\
\textbf{Formulação equivalente:} $\oint\vec E\cdot\mathrm{d}\vec\ell=0$ em
qualquer curva fechada. O campo eletrostático é \textbf{conservativo}, e essa é
exatamente a lei das malhas de Kirchhoff em sua forma de campo.\\
\textbf{Ressalva importante:} isso vale para campos \emph{eletrostáticos}.
Campos elétricos induzidos por variação de fluxo magnético \emph{não} são
conservativos, e para eles o potencial escalar deixa de existir --- assunto da
lei de Faraday.}
\end{questao}

\begin{questao}[4]{}
\textbf{(Projeto --- $V=0$ com $E\neq0$.)} Projete um par de cargas de sinais
opostos e módulos \emph{diferentes} tal que exista um ponto do segmento onde
$V=0$ mas $\vec E\neq\vec0$.
\begin{itensq}
\item Determine a condição geral.
\item Escolha valores e localize o ponto.
\item Verifique que o campo não é nulo ali.
\end{itensq}
\resposta{Condição: $|q_1|/r_1=|q_2|/r_2$; com $+4\,\si{\micro\coulomb}$ e
$-1\,\si{\micro\coulomb}$ a $\SI{0.5}{\meter}$, o ponto é
$x=\SI{0.4}{\meter}$}
\resolucao{(a) \textbf{Condição para $V=0$:}
\[
\frac{\ke q_1}{r_1}+\frac{\ke q_2}{r_2}=0
\;\Longrightarrow\;
\frac{|q_1|}{r_1}=\frac{|q_2|}{r_2},
\]
com $q_1$ e $q_2$ de sinais \emph{opostos}.\\
\textbf{Condição para $\vec E=\vec0$ no mesmo ponto:}
$\dfrac{|q_1|}{r_1^{2}}=\dfrac{|q_2|}{r_2^{2}}$.\\
As duas só coincidem se $r_1=r_2$ --- e então $|q_1|=|q_2|$, que é o caso
excluído pelo enunciado. \textbf{Com módulos diferentes, é impossível anular
ambos no mesmo ponto.}\\
(b) Tome $q_1=+\SI{4}{\micro\coulomb}$ em $x=0$ e
$q_2=-\SI{1}{\micro\coulomb}$ em $x=\SI{0.5}{\meter}$. De
$\dfrac{4}{x}=\dfrac{1}{0{,}5-x}$ vem $x=\SI{0.4}{\meter}$.\\
\emph{(Confira: $r_1=0{,}4$ e $r_2=0{,}1$, com
$4/0{,}4=1/0{,}1=10$ \checkmark.)}\\
(c) \textbf{Campo no ponto.} Ambos apontam no sentido $+x$ (afastando-se de
$q_1$, aproximando-se de $q_2$) e portanto se somam:
\[
E=\frac{\pot{9}{9}(\pot{4}{-6})}{0{,}16}
+\frac{\pot{9}{9}(\pot{1}{-6})}{0{,}01}
=\pot{2.25}{5}+\pot{9}{5}=\pot{1.125}{6}\,\si{\volt\per\meter}.
\]
Mais de um megavolt por metro --- longe de nulo.\\
\textbf{Interpretação:} uma carga colocada ali teria energia potencial nula, mas
sofreria força intensa. \textbf{$V$ e $\vec E$ carregam informações
diferentes:} $V$ é o valor, $\vec E$ é a taxa de variação.}
\end{questao}

\begin{questao}[4]{}
\textbf{(Estabilidade em um mínimo de potencial.)} Uma curva
$V(x)$ apresenta mínimo local em $x_0$.
\begin{itensq}
\item Analise a estabilidade de uma carga positiva colocada em $x_0$.
\item Repita para uma carga negativa.
\item Relacione com o teorema de Earnshaw.
\end{itensq}
\resposta{Positiva: \textbf{instável}; negativa: estável em $x$ --- mas nenhuma
é estável em três dimensões}
\resolucao{(a) A energia potencial de uma carga positiva é $U=qV$, com
$q>0$ --- ela tem a \emph{mesma} forma de $V$ e, portanto, também um mínimo em
$x_0$. Mínimo de energia significa \textbf{equilíbrio estável}... \emph{ao longo
de $x$}.\\
\emph{(Atenção ao raciocínio: é o mínimo de $U$, e não o de $V$, que define
estabilidade. Para carga positiva os dois coincidem; para negativa, invertem-se.)}\\
(b) Para carga negativa, $U=qV$ com $q<0$ \textbf{inverte} a curva: onde $V$
tem mínimo, $U$ tem \emph{máximo}. Equilíbrio \textbf{instável} em $x$.\\
(c) \textbf{Teorema de Earnshaw.} Em uma região sem cargas, o potencial
eletrostático satisfaz a equação de Laplace, $\nabla^{2}V=0$, ou seja
\[
\frac{\partial^{2}V}{\partial x^{2}}
+\frac{\partial^{2}V}{\partial y^{2}}
+\frac{\partial^{2}V}{\partial z^{2}}=0.
\]
Se $V$ tivesse um mínimo \emph{tridimensional} em $x_0$, as três segundas
derivadas seriam positivas e a soma não poderia ser zero. \textbf{Contradição.}\\
Logo, um mínimo em uma direção obriga a haver \emph{máximo} em outra --- todo
ponto crítico é um \textbf{ponto de sela}.\\
\textbf{Conclusão:} nenhuma carga pode ser mantida em equilíbrio estável apenas
por forças eletrostáticas, seja ela positiva ou negativa. A análise
unidimensional do item (a) é enganosa justamente por ignorar as outras duas
direções.\\
\textbf{Como se contorna na prática:} campos variáveis no tempo (armadilha de
Paul), campos magnéticos, ou realimentação ativa --- todas soluções que escapam
da hipótese eletrostática.}
\end{questao}

\begin{questao}[4]{}
\textbf{(Condutor em equilíbrio.)} Demonstre que toda a superfície de um
condutor em equilíbrio eletrostático é equipotencial.
\begin{itensq}
\item Prove a partir da condição de equilíbrio.
\item Mostre que o interior também está no mesmo potencial.
\item Deduza a orientação do campo na superfície.
\end{itensq}
\resposta{Campo interno nulo $\Rightarrow$ $V$ constante em todo o condutor; o
campo externo é \textbf{perpendicular} à superfície}
\resolucao{(a) \textbf{Condição de equilíbrio:} não há movimento de cargas
livres, logo o campo dentro do metal é \textbf{nulo}. Se houvesse campo, as
cargas se moveriam e o equilíbrio não existiria.\\
Tomando dois pontos $A$ e $B$ quaisquer do condutor e um caminho
\emph{interno} entre eles:
\[
V_A-V_B=\int_A^B\vec E\cdot\mathrm{d}\vec\ell=\int_A^B\vec 0\cdot\mathrm{d}\vec\ell=0.
\]
Portanto $V_A=V_B$ --- \textbf{todo o condutor é equipotencial}.\\
(b) O argumento não distinguiu superfície de interior: ele vale para
\emph{qualquer} par de pontos ligados por caminho dentro do metal. Superfície e
interior estão no mesmo potencial, e o condutor inteiro é um volume
equipotencial.\\
(c) \textbf{Orientação do campo externo.} Se houvesse componente
\emph{tangencial} à superfície, ela empurraria as cargas livres ao longo dela --
contradizendo o equilíbrio. Logo o campo externo é \textbf{perpendicular} à
superfície em todo ponto.\\
\textbf{Consequências, todas verificáveis:}
\begin{itemize}
\item A carga em excesso reside \emph{integralmente} na superfície externa.
\item A densidade superficial é maior onde a curvatura é maior (poder das
pontas).
\item Uma cavidade vazia dentro do condutor tem campo nulo --- é a blindagem
eletrostática.
\item Ligar dois condutores por fio iguala seus potenciais, e a carga se
redistribui proporcionalmente às capacitâncias.
\end{itemize}}
\end{questao}

\begin{questao}[4]{}
\textbf{(Do potencial ao campo.)} O potencial em uma região é
$V(x,y)=\SI{100}{}-\SI{20}{}\,x+\SI{5}{}\,y^{2}$ (volts, com $x,y$ em metros).
\begin{itensq}
\item Determine $\vec E(x,y)$.
\item Avalie em $(x,y)=(2;\,1)$.
\item Determine a forma das equipotenciais.
\end{itensq}
\resposta{$\vec E=(20;\,-10y)\,\si{\volt\per\meter}$; em $(2;1)$,
$\vec E=(20;\,-10)$ com $|E|=\SI{22.4}{\volt\per\meter}$}
\resolucao{(a) $\vec E=-\nabla V$, componente a componente:
\[
E_x=-\frac{\partial V}{\partial x}=-(-20)=\SI{+20}{\volt\per\meter};
\qquad
E_y=-\frac{\partial V}{\partial y}=-(10y)=-10y.
\]
Note que $E_x$ é \emph{uniforme}, enquanto $E_y$ cresce linearmente com $y$.\\
(b) Em $(2;\,1)$:
\[
\vec E=(20;\,-10)\,\si{\volt\per\meter};
\qquad
|E|=\sqrt{400+100}=\SI{22.4}{\volt\per\meter},
\]
formando $\ang{-26.6}$ com o eixo $x$.\\
\textbf{Observe:} o valor de $V$ no ponto ($V=100-40+5=\SI{65}{\volt}$) não
entrou no cálculo do campo. Só as \emph{derivadas} importam --- somar uma
constante a $V$ não altera $\vec E$.\\
(c) \textbf{Equipotenciais:} $100-20x+5y^{2}=\text{const}$, ou seja
\[
x=\frac{5y^{2}+100-C}{20},
\]
uma família de \textbf{parábolas} com eixo horizontal.\\
\textbf{Verificação de coerência:} o vetor $\vec E$ calculado deve ser
perpendicular à parábola que passa por $(2;1)$. A tangente à curva ali tem
inclinação $\mathrm{d}y/\mathrm{d}x=2$, e o campo tem inclinação
$-10/20=-1/2$ --- produto $-1$, portanto perpendiculares \checkmark.}
\end{questao}

\begin{questao}[4]{}
\textbf{(Equipotenciais e linhas de campo.)} Compare os dois conjuntos de
curvas para duas cargas positivas iguais.
\begin{itensq}
\item Descreva as equipotenciais perto de cada carga e muito longe do par.
\item Descreva as linhas de campo, incluindo o ponto médio.
\item Estabeleça as regras gerais que relacionam os dois conjuntos.
\end{itensq}
\resposta{Perto: esferas em torno de cada carga; longe: esferas em torno do par;
o ponto médio é ponto de sela do potencial}
\resolucao{(a) \textbf{Equipotenciais.}
\begin{itemize}
\item \emph{Muito perto de uma carga:} a contribuição dela domina, e as
superfícies são praticamente \textbf{esferas} centradas nessa carga.
\item \emph{Muito longe do par:} o sistema parece uma única carga $2Q$, e as
superfícies voltam a ser esferas, agora centradas no \textbf{centro do par}.
\item \emph{Na região intermediária:} as superfícies deixam de ser esferas e
assumem forma de "amendoim", chegando a se fundir em uma única superfície que
envolve as duas cargas.
\end{itemize}
(b) \textbf{Linhas de campo.} Saem radialmente de cada carga, curvam-se
afastando-se uma da outra e seguem para o infinito. \textbf{No ponto médio o
campo é nulo} --- as linhas não passam por ali. Não há linha ligando as duas
cargas, pois ambas são fontes.\\
(c) \textbf{Regras gerais.}
\begin{enumerate}
\item As linhas de campo são \textbf{perpendiculares} às equipotenciais em todo
ponto.
\item O campo aponta no sentido de $V$ \textbf{decrescente}.
\item Onde as equipotenciais estão \textbf{mais próximas}, o campo é mais
\textbf{intenso} --- pois $|E|=|\mathrm{d}V/\mathrm{d}n|$.
\item Equipotenciais nunca se cruzam (um ponto teria dois potenciais); linhas de
campo também não (haveria duas direções de força).
\item Onde $\vec E=\vec0$, as duas regras de perpendicularidade perdem sentido:
é um \textbf{ponto de sela} do potencial, e ali as equipotenciais podem se
cruzar.
\end{enumerate}
\textbf{O ponto médio é justamente esse caso excepcional:} $V$ tem mínimo ao
longo da linha que une as cargas e máximo na direção perpendicular --- exatamente
a estrutura de sela exigida pelo teorema de Earnshaw.}
\end{questao}

% END BANCO COMPLEMENTAR

\gabarito[1.4 Potencial elétrico e superfícies equipotenciais]
